%% capitulos/1-introducao.tex

\chapter{Introdução}
\label{cap:introducao}

Este capítulo apresenta a contextualização do problema de pesquisa,
a motivação, os objetivos, a metodologia adotada e a estrutura
desta dissertação.

% ── Dica: use \section, \subsection etc. para subdivisões

\section{Contextualização e Motivação}
\label{sec:contextualizacao}

Descreva aqui o contexto da área e o que motiva a pesquisa.
Cite autores relevantes \cite{exemplo2023}.

\section{Problema de Pesquisa}
\label{sec:problema}

Enuncie o problema de pesquisa de forma clara e objetiva.

\section{Objetivos}
\label{sec:objetivos}

\subsection{Objetivo Geral}

O objetivo geral desta dissertação é \ldots

\subsection{Objetivos Específicos}

Para atingir o objetivo geral, os seguintes objetivos específicos
foram definidos:

\begin{enumerate}
  \item Objetivo específico 1;
  \item Objetivo específico 2;
  \item Objetivo específico 3.
\end{enumerate}

\section{Metodologia}
\label{sec:metodologia-intro}

% Breve descrição da abordagem metodológica (detalhada no cap. 4)
Esta pesquisa adota o paradigma da Design Science Research
\cite{Hevner2004}, \ldots

\section{Contribuições}
\label{sec:contribuicoes}

As principais contribuições deste trabalho são:

\begin{itemize}
  \item Contribuição 1;
  \item Contribuição 2.
\end{itemize}

\section{Estrutura da Dissertação}
\label{sec:estrutura}

Este trabalho está organizado da seguinte forma:
o Capítulo~\ref{cap:fundamentacao} apresenta a fundamentação teórica;
% Adicione os demais capítulos aqui
